\documentclass[11pt,a4paper]{article}

% --- 1. Encoding, Fonts, Language ---
\usepackage[utf8]{inputenc}
\usepackage[T1]{fontenc}
\usepackage{lmodern}
\usepackage[english]{babel}

% --- 2. Mathematics ---
\usepackage{amsmath, amssymb, amsthm}
\usepackage{empheq}

% --- Theorem environments (manual numbering kept consistent with cross-references) ---
\theoremstyle{plain}
\newtheorem*{thmA}{Theorem 1 (Exclusion of the static lattice)}
\newtheorem*{thmB}{Theorem 2 (Exclusion of the self-bound liquid)}
\newtheorem*{thmC}{Theorem 3 (Exact covariance of the scalar sector)}
\newtheorem*{thmD}{Theorem 4 (Uniqueness of the Johns node from substrate constraints)}
\theoremstyle{remark}
\newtheorem*{remarkstar}{Remark}

% --- 3. Geometry and Layout ---
\usepackage[margin=2.5cm]{geometry}
\usepackage{titlesec}
\usepackage{enumitem}

% --- 4. Tables, Figures, Captions ---
\usepackage{graphicx}
\graphicspath{{../figures/}}
\usepackage{booktabs}
\usepackage{longtable}
\usepackage{array}
\usepackage{caption}

% --- 5. TikZ (2D & 3D) ---
\usepackage{tikz}
\usepackage{tikz-3dplot}
\usetikzlibrary{arrows.meta, backgrounds, shadings, positioning, calc}

\tikzset{
    picto_box/.style={
        draw=gray!60,
        thick,
        minimum width=1.5cm,
        minimum height=1.5cm,
        inner sep=5pt,
        rounded corners=2pt
    },
    pole/.style={circle, fill=black, inner sep=1.2pt}
}

% --- 6. Hyperlinks (always load last!) ---
\usepackage{hyperref}
\hypersetup{
    colorlinks=true,
    linkcolor=blue,
    urlcolor=cyan,
    pdftitle={The Matter Sector of the Dipolar String Model},
}

% --- 7. Document metadata ---
\title{\textbf{The Matter Sector of the Dipolar String Model:}\\ \large \textsc{a redraw from the author's constitutive statements, with its exclusions}}
\author{R. St-Pierre\thanks{Independent researcher, Qu\'ebec, Canada. This research received no external funding. The author declares no conflict of interest.}}
\date{September 2026 --- admin@dipolarstrings.org --- draft, not deposited}

% --- 8. Custom macros ---
\newcommand{\dqd}{\ensuremath{\ooalign{$\bigcirc$\cr\hfil\scalebox{0.7}{$=$}\hfil}}}
\newcommand{\Ncd}{\ensuremath{N_{\mathrm{DS}}}}
\newcommand{\lb}{\ensuremath{\bar\lambda_e}}
\newcommand{\lone}{\ensuremath{\ell_1}}
\newcommand{\R}[1]{R#1}

% Every claim carries its maturity level.  Six levels, as in the main manuscript:
% Derived, Conditionally derived, Calibrated, Postulated, Reparametrized, Closed.
% A ratio obtained by fitting a measured mass is Reparametrized, not Derived.
\newcommand{\status}[1]{\medskip\noindent\textbf{Status: #1.}\ }
\setlength{\emergencystretch}{2.5em}

\begin{document}

\maketitle

\begin{abstract}
\noindent\textbf{Background.} The Dipolar String (DS) framework is a theory of the medium. An
impedance-matched, purely electromagnetic vacuum has no intrinsic scale, so the electron mass is
an input, and the matter sector was out of scope. Between 14 and 19 September 2026 the sector was
redrawn from eight constitutive statements of the author (R1--R8), computing at each step what the
statement imposes and what it breaks; the record is 103 steps (R1--R103) backed by 98 scripts.

\noindent\textbf{Methods.} Each statement became a self-contained script with PASS/FAIL checks and
a verdict (derived, conditional, coincidence, identity, excluded), recorded in
\texttt{matter/redraw/BASE.md}. Inputs: $\alpha$, $m_e$, the exponent $2\pi$ (pinned by the muon
and tau masses to $0.15\,\%$), the nucleon strand count 9, and the section of the ribbon. The
computation was performed by a second AI system working from the author's statements; it is not
independent verification (see the Disclosure).

\noindent\textbf{Results.} Three strands carrying $-1/3$, $0$, $+1/3$ give exactly the observed
charge spectrum and a colour label. A circulation quantum $\pi\hbar c$ per closed path, with the
path scale $\lone(n) = (2\pi\lb/3)(3/n)^{2\pi}$, gives the quark as an electron-type circuit at
$\lone(9) = 0.81$~fm (254~MeV), the nucleon as three static circuits (763~MeV) plus a spin-carrying
ring, the string tension $\sqrt\sigma = 430$~MeV (lattice 420--440), the neutron--proton splitting
1.27--1.34~MeV (measured 1.293) and the $\Delta$--$N$ gap within $15\,\%$. The lepton masses are the
Koide form of a $C_3$ mass-root operator whose $45^\circ$ is the half-and-half partition of
$g = 2$; the phase $2/9$ is posed. The lifetime law $\tau = T\,(M/m)^4/C(n)$ reproduces the muon and
the tau, with $M = (96\pi^2)^{1/4}/\sqrt{G_F}$ an identity, not a derivation. The electron itself
does not survive: any charge carried by a 386~fm ring has the form factor $j_0(qR)$, excluded by
Compton, $e$--$p$ and LEP data; a Bohr magneton carried by a charge loop needs $R \ge \lb$ at
$v \le c$; the ring's numbers are Compton identities fitted to Dirac. Decision: the
electromagnetic electron is Dirac's.

\noindent\textbf{Conclusions.} What the framework can now claim is a set of hadronic and
generational ratios given one anchor, each at its stated level. What it cannot claim is an
extended electron, a mechanism for the phase $2/9$, the weak scale, fermionic statistics, or a
local account of Bell correlations.
\end{abstract}

%==============================================================================
\section{Scope, and what this paper does not claim}
%==============================================================================

The main manuscript \cite{dsmodel} is a theory of the medium. It withdrew its muon, tau, $W$,
$Z$ and $H$ table entries as reverse-engineered mass fits without predictive content, and it
establishes that a purely electromagnetic, impedance-matched vacuum is conformally invariant:
there is no intrinsic scale, $m_e$ is structurally an input, and what is attainable is a spectrum
of \emph{ratios} given the single anchor $\lb = \hbar/m_e c_0 = 386.16$~fm. This paper accepts
that boundary and attempts one sector of matter inside it.

Four things are stated before any result.

\begin{enumerate}[nosep]
\item Every number below is a ratio to $m_e$ or a length in units of $\lb$. No absolute mass is
claimed. Where an index was chosen to reproduce the mass it explains, the entry is labelled
\emph{Reparametrized} and is not counted as a result.
\item The inputs are five: $\alpha$, $m_e$, the exponent $2\pi$ of the path-scale ladder, the
strand count 9 of the nucleon, and the square section of the ribbon. The exponent and the count
are structural inputs that are not derived; this is stated wherever a result depends on them.
\item The electron is \emph{not} a result of this paper. Steps R99--R103 show that the electron
representation the redraw had built, a closed ring of radius $\lb$ carrying the charge $e$, is
excluded by the measured electromagnetic form factors, and that its numbers ($S = \hbar/2$,
$\mu_B$, $g = 2$) are identities of the choice $R = \lb$, $E_{\rm circ} = m/2$, $q = e$ rather than
outputs. The decision taken (R103) is that the electromagnetic electron is Dirac's point particle.
The sector's claims are therefore hadronic and generational.
\item The computation was done by a second AI system on the author's statements, with the
repository's constraints in view. This is a shared prior, not independent replication
(Section~\ref{sec:disclosure}).
\end{enumerate}

%==============================================================================
\section{Inherited framework}\label{sec:inherited}
%==============================================================================

From the manuscript: the anchor $R_3 = \hbar/m_e c_0 = \lb$, hence $\lone = 2\pi R_3/3 = 808.8$~fm
for the electron's three strings; the vacuum is a weave of neutral dipolar strings (DQD, \dqd),
bifilar lines matched to the impedance of free space $Z_0$, with the matched pair separation
$D/r = 2\cosh\pi$; the states in scope were the photon ($\Ncd = 1$), the DQD ($\Ncd = 2$) and the
electron ($\Ncd = 3$). The electron impedance $Z_e \approx 0.73\,Z_0$ was filed \emph{Calibrated}
because it depends on the undetermined aspect ratio $R_3/r \approx 37.1$.

Three of these are modified by the redraw, each by computation and each recorded.

\begin{itemize}[nosep]
\item \textbf{The mass law.} ``Mass is inductance'' is not used. The repository's own audit
(\texttt{CONSTRAINTS.md}~\S1b--1c) had shown that the inductive self-energy of the electron loop is
$0.4\,\%$ of $m_e c^2$ and that no coil of 27 strings reaches the proton. The redraw's mass is the
energy of a circulation quantum per closed path, $E = \pi\hbar c/\ell_{\rm path}$ (R32, R40, R47),
with the path scale set by a ladder in the strand count (Section~\ref{sec:ladder}).
\item \textbf{The tube radius.} The manuscript's fixed $r = R_3/37.1$ capped a mode energy at
$m/m_e \le 37$. The redraw finds the brick has no size: by the two-dimensional method of moments,
the capacitance per length of two square conductors depends on the aspect ratio only (identical at
$\times 100$ to $10^{-9}$), so $Z$ and the energy per length of a bifilar line are scale
invariant and the section follows the circuit, $w/\lone = 6/\pi^3$ at every rung (R71,
\texttt{brick\_scale.py}). The cap does not apply; matching fixes the \emph{shape} of the section
(square, R17), not its size.
\item \textbf{$Z_e = 0.73\,Z_0$.} Read as the impedance of the electron's circuit it would leak
$2.4\,\%$ per turn and the electron would live $3\times10^{-19}$~s (R65). The electron's circuit is
matched to $Z_0$ to better than $10^{-28}$; the calibrated number is not the circuit impedance.
\end{itemize}

The repository's other constraints are respected as written: charge is not carried by the
connectivity of the network but by an unpaired branch of the $\pm e/3$ fluid (\S2, and R28
below); no Dzyaloshinskii--Moriya stabilization is invoked (\S3); neutrino masses are not
attempted (\S7).

%==============================================================================
\section{The obstruction, and what remains attainable}\label{sec:obstruction}
%==============================================================================

Four independent analyses in the manuscript converge on the absence of an intrinsic scale: the
equation of state $P = u/3$, the single-velocity character of the linear lattice, the pure
ultraviolet character of the zero-point sum over the braided defect, and the fact that a
transmission line's soliton-to-quantum mass ratio is fixed by its impedance alone, so that a
matched line carries the coupling $2\pi\alpha$ and every soliton weighs $O(1/\alpha)$ quanta.
Coiling a line does not raise the coupling. The corollary governs everything below: ratios yes,
absolute masses no, and no confinement-type regime reachable by winding.

The redraw confirms the corollary from inside. Winding a charged strand into a solenoid does
amplify the energy per length quadratically in the turns per core radius (R33), but the force
along the axis never exceeds the energy per unit \emph{path}, $dE/dL = u_{\rm path}/\gamma \le
u_{\rm path}$ (R39): the winding route to the string tension closes, and what remains is the
energy per path of the matched line itself, which is where the tension of
Section~\ref{sec:tension} comes from. Likewise a fixed linear charge density is excluded by the
electron (its strands carry $4\times10^{-4}\,e$/fm against the $7\,e$/fm a lattice-strength tension
would need, R31), and a flux tube is not a mechanism of the base as written, since a conducting
fluid screens a pole instead of channelling it (R30).

What is attainable, and what this paper reports, is the following: a charge spectrum; a ladder of
path scales in the strand count; hadronic ratios (quark, nucleon, $\Delta$, $n$--$p$, tension,
nuclear force) given the count 9; the $C_3$ structure of the lepton masses; a lifetime law; and a
set of exclusions, the largest of which is the extended electron.

%==============================================================================
\section{Results}\label{sec:results}
%==============================================================================

Each subsection names the steps of \texttt{matter/redraw/BASE.md} and the scripts of
\texttt{matter/redraw/} that produce its numbers; every script prints a \texttt{RESULT: n/m PASS}
line and exits non-zero on any failure. Verdicts in the record are \emph{derived},
\emph{conditional}, \emph{coincidence}, \emph{identity} and \emph{excluded}; they are mapped to
the six levels of the manuscript at the end of each subsection.

%------------------------------------------------------------------------------
\subsection{Charge: three strands and the unpaired branch}\label{sec:charge}
%------------------------------------------------------------------------------

The author's statements R1--R8 place the charge in the polarized fluid of the string, with a pole
at each end. The repository's note of 15 September had already established the divergence source:
a complete DQD is neutral, an unpaired branch is a free charge $e/3$. Chains of charged strings,
with $q = (n_+ - n_-)e/3$, allow charges $\pm 4/3$ and $\pm 5/3$ that are never observed (R18), and
a star of three charged strings cannot carry $+2/3$ (R22). The repair is the \emph{neutral
strand} (R28, \texttt{three\_strand.py}): every fermion is three strands, each carrying $-1/3$,
$0$ or $+1/3$:
\[
e^- = (-,-,-),\quad \nu = (0,0,0),\quad u = (+,+,0),\quad d = (-,0,0),
\]
antiparticles by sign reversal. This gives the observed spectrum $\{0, \pm1/3, \pm2/3, \pm1\}$
exactly and nothing else; colour is the position of the odd strand (three arrangements for $u$ and
$d$, one for $e$ and $\nu$; a baryon is colourless when its three odd strands occupy three
different positions); $p = uud$ and $n = udd$ both have nine strands, which is what
$m_p \approx m_n$ required; the pions have six. A neutral strand is a complete DQD, the unit of the
medium. The rule ``charged strands carry the mode, the total count sets the scale'' keeps
$m_\nu = 0$ (KATRIN bound $0.45$~eV \cite{katrin}) with $e$, $\mu$, $\tau$ unchanged; it is posed.

The same assignment is that of the three-strand braid model of Bilson-Thompson
\cite{bilsonthompson}; the redraw reaches it from the charge spectrum alone and does not adopt the
braid dynamics.

\status{Derived, given the postulated strand charges $\{-1/3, 0, +1/3\}$; the neutrino rule is
Postulated}

%------------------------------------------------------------------------------
\subsection{The circulation quantum and the ladder of path scales}\label{sec:ladder}
%------------------------------------------------------------------------------

\paragraph{The fluid is a vortex, not a wave.} A standing wave on a guide has harmonics and a
transverse cut-off; for the electron's ribbon ($w = 4\lb/\pi^2 = 156.5$~fm) the transverse mode
sits at $\hbar c\pi/w = 3.96$~MeV, and the arc harmonics at 0.77 and 1.02~MeV, none observed (R17,
R47). Read as a stationary circulation carrying one conserved quantum per closed circuit, the
fluid has neither: the energy of the circuit is $E = \pi\hbar c/\ell_{\rm path}$, and the identity
$4\pi K q^2 = \pi\hbar c$ with $K = \alpha\hbar c$ and $q = e/(2\sqrt\alpha) = 5.85\,e$ makes it
equal to the half-wave energy for every path, so every number obtained in the wave reading is
kept (R32, R40, R47; \texttt{vortex.py}, \texttt{circuit\_rule.py}, \texttt{transverse\_mode.py}).
The ``vortex charge'' $q$ is not a charge that circulates (R53); it is the coefficient of this
identity. The quantum $\pi\hbar c$ per path is an action of half a turn, $\oint p\,dl = \pi\hbar$
(R86), and can be obtained as half the whole quantum $h$ of a periodic neutral mother loop broken
symmetrically into two locked daughters (R87--R88), under three named inputs (one quantum on the
mother, momentum conserved at the break, C-symmetric sharing).

\paragraph{The ladder.} With the mode on the three charged strands of a lepton, the path is
$3\lone(n)$ and the masses require $\lone(n) = \lone(3)(3/n)^{2\pi}$ for $n = 3, 7, 11$
(\texttt{lepton\_ladder.py}, \texttt{exponent\_2pi.py}). The exponent is pinned by the data:
$(e,\mu)$ alone give $6.2926$, $(e,\tau)$ alone $6.2758$, and $2\pi = 6.2832$ lies at $0.016\,\%$
of their midpoint; among about a thousand simple two-factor constants it is the only one within
$0.15\,\%$ (R55). It is \emph{not} derived, and it is not the exact law: the ladder gives the muon
at $-0.8\,\%$ and the tau at $+1.0\,\%$, while the measured masses obey Koide's relation to
$10^{-5}$; the exponent that would make the ladder exactly Koide, $6.25$, is outside the data
window. Any added parameter makes the ladder worse (R73). The ladder is therefore an
approximation at the percent level, the exact structure being that of Section~\ref{sec:koide}.
The indices $3, 7, 11$ follow from the neutral strand: a generation adds four neutral strands, two
DQD (R28); a single DQD, of spin 1, cannot be added to a spin-$\tfrac12$ ring and keep its spin,
two can (R70), so $n = 5$ (12.7~MeV, absent) is forbidden. The next rung $n = 15$ would be a
12.6~GeV charged lepton living $10^{-17}$~s, below the LEP bound of 100.8~GeV \cite{pdg}: the
ladder \emph{as a mass law} is excluded by the absence of a fourth generation (R76), which is what
forces the reading above.

\status{Postulated (one quantum per circuit; R87 conditionally derives it); Reparametrized for the
exponent $2\pi$; Derived for the step of four neutral strands and the exclusion of $n = 5$}

%------------------------------------------------------------------------------
\subsection{Leptons: the Koide form as a $C_3$ mass-root operator}\label{sec:koide}
%------------------------------------------------------------------------------

Koide's relation \cite{koide}, $Q = \sum m/(\sum\sqrt m)^2 = 2/3$, holds for $(e, \mu, \tau)$ to
$10^{-5}$ and means that the vector $(\sqrt{m_e}, \sqrt{m_\mu}, \sqrt{m_\tau})$ makes exactly
$45^\circ$ with the diagonal $(1,1,1)$ \cite{foot}; with the Brannen phase $2/9$ \cite{brannen} it
reproduces $m_\mu$ from $m_e$ to $+0.001\,\%$ and $m_\tau$ to $+0.007\,\%$ (R60,
\texttt{generation\_holonomy.py}).

The redraw's structure (R76--R81; \texttt{c3\_closure.py}, \texttt{koide\_hamiltonian.py},
\texttt{koide\_from\_g2.py}):

\begin{itemize}[nosep]
\item Adding two DQD is the cyclic shift $U$ of the three strand positions ($4 \equiv 1 \bmod 3$);
$U^3 = I$, $U \ne I$, $U^2 \ne I$, so a $Z_3$-symmetric hermitian \emph{mass-root} operator
$C = aI + bU + \bar b U^2$ (eigenvalues $\lambda_k = \sqrt{m_k}$) has exactly three eigenvalues
$a + 2|b|\cos(\varphi + 2\pi k/3)$, the Koide--Brannen form, and the fourth application is the
first. Koide's $45^\circ$ is $a^2 = 2|b|^2$.
\item For any hermitian $3\times3$ operator, $Q = 1/3 + (\|C_{\rm dev}\|/\|C_{\rm iso}\|)^2/3$;
Koide is the equality of the isotropic and deviatoric norms (measured ratio $0.999991$). The
equality is not selected by any extremum (the isotropic fraction decreases monotonically from
equal masses to rank one); it must be a conservation, one quantum per irreducible channel, not per
degree of freedom (which would give $Q = 1$ and two massless leptons).
\item For a circulant, $\|C_{\rm iso}\|^2$ is the on-site weight and $\|C_{\rm dev}\|^2$ the
hopping weight, and ${\rm tr}\,C^2$ is the family's total mass. If the mass-root operator is the
amplitude matrix of the three-position network (junction = node, circulation = link), then the
half-static, half-circulating partition that gives $g = 2$ for each lepton (R53--R54), summed
over the family, gives $a^2 = E_{\rm stat}/3 = 313.84$~MeV and $|b|^2 = E_{\rm circ}/6$, i.e.
$a^2 = 2|b|^2$ without a $\sqrt2$ put by hand ($a^2$ from Koide $+2/9+m_e$: 313.86~MeV). The
pairing holds at the family level only; it does not touch $\varphi$.
\item The phase $\varphi = 2/9$ is posed. The base can write it exactly as the energy of one
junction ($m_e c^2/6$) during one radian of the fluid's orbit at the radius of the three-quarter
turn, or as the product of charges $|q_u q_d|$, but has no rule making either the phase of the
doublet (R58--R63). A Berry-phase route gives $a\cdot 2\pi/3 = 0.39$~rad, excluded.
\item The topological lock that binds a pair of DQD is real ($Lk = Tw + Wr$ conserved with the
frame locked to the medium, writhe $1/3$ at $r_0 = 0.1785$, R61, R77), but it binds a third pair
too (writhe $1$ at $r_0 = 0.376$): the absence of a fourth generation comes from the spectrum of
$C$, not from the lock. $N_\nu = 2.9963 \pm 0.0074$ \cite{pdg} is the three-sector count.
\end{itemize}

\status{Conditionally derived for the form and the $45^\circ$ (on the identification of the
mass-root operator with the network amplitude matrix); Postulated for $\varphi = 2/9$;
Derived that the ladder cannot be corrected into Koide}

%------------------------------------------------------------------------------
\subsection{Quarks and nucleons}\label{sec:nucleon}
%------------------------------------------------------------------------------

Under the rule of one quantum per closed circuit, a circuit of $k$ strands in series at the
nucleon rung has the energy $\pi\hbar c/(k\,\lone(9))$ with $\lone(9) = 0.8126$~fm and
$\pi\hbar c/\lone(9) = 763$~MeV. The static part of the proton required by its magnetic moment
(770--813~MeV, R25) must be the sum over its circuits; of the 30 partitions of nine strands, only
$\{3,3,3\}$ has $\sum 1/k_i = 1$ (R41, \texttt{circuit\_partition.py}): \textbf{three circuits of
three strands, one per quark, each the electron's object at the scale $\lone(9)$}, 254~MeV each,
the energy ratio to the electron being $3^{2\pi} = 996$ exactly. The remainder is a single ring
that carries the spin and the moment (R72, \texttt{nucleon\_single\_count.py}); with
$S = R\,E_{\rm circ}/c = \hbar/2$ its radius is $0.56$--$0.59$~fm from $\mu_p$, and the
$\Delta$--$N$ gap, read as the same ring with three times the circulation ($3\hbar/2$), is
$2E_{\rm circ}(N) = 250$--$336$~MeV against 294 measured (R48, \texttt{delta\_shape.py}).
Inverting, $E_{\rm circ}(N) = 147$~MeV at $R = 0.671$~fm gives $N = 910$~MeV ($-3.0\,\%$) and
$\Delta = 1204$~MeV ($-2.3\,\%$) with one radius; that radius equals $\sqrt3$ times the circuit
radius $3\lone/2\pi$ to $0.1\,\%$, and the envelope of three tangent circuits is $0.836$~fm, the
proton charge radius to $0.6\,\%$, both coincidences without a selecting mechanism (R49).

\begin{table}[h]
\centering\small
\begin{tabular}{lccl}
\toprule
quantity & this work & measured / lattice & inputs and level \\
\midrule
quark circuit & 254 MeV & constituent $\sim 313$ & $2\pi$, 9; Conditionally derived \\
nucleon, static part & 763 MeV & --- & partition $\{3,3,3\}$; Derived given the rule \\
nucleon, total & 910--938 MeV & 938.3 & ring radius posed; Conditionally derived \\
$\Delta - N$ & 250--336 MeV & 294 & R25 radii; Conditionally derived \\
$m_n - m_p$, strong part & 2.47 MeV & $2.52 \pm 0.29$ \cite{bmw} & $(2\alpha/3)\,3^{2\pi} m_e c^2$; Derived \\
$m_n - m_p$, total & 1.27--1.34 MeV & 1.293 & section, pole shape; Conditionally derived \\
$\mu_p/\mu_n$ & $-3/2$ & $-1.460$ & sharing $2/3, 2/3, -1/3$; Postulated \\
$\sqrt{\langle r^2\rangle_p}$ & 0.81--0.84 fm & 0.841 & coincidence \\
\bottomrule
\end{tabular}
\caption{Hadronic ratios. All numbers hang on $\lb$, $\alpha$, the exponent $2\pi$ and the count 9.}
\end{table}

The neutron--proton splitting (R42--R45; \texttt{pn\_splitting.py}, \texttt{coulomb\_geometry.py},
\texttt{log\_constant.py}, \texttt{partner\_screening.py}): $p = uud$ and $n = udd$ are the same
three circuits and the same junction graph, so the 763~MeV and the joints cancel; the neutron has
one more neutral strand, a complete DQD whose bifilar energy at the nucleon rung is
$(e/3)^2/(\epsilon_0\lone(9)) = 2.47$~MeV, the strong part of the splitting (lattice
$2.52 \pm 0.29$~MeV \cite{bmw}). The proton's Coulomb energy, with the charge on the charged strands
(two for $u$, one for $d$), the star geometry (mutual term zero for $p$, $-0.34$~MeV for $n$) and
the ribbon section as the transverse cut-off (a thin strip of width $w_9 = 0.157$~fm has the
capacitance of a cylinder of radius $w/4$), gives $m_n - m_p = 1.32$~MeV ($+2.4\,\%$), and
$1.27$~MeV ($-1.8\,\%$) with the lowest-energy pole shape in the square section (R72); screening by
the neutral neighbours changes it by less than $1.5\,\%$. The section enters only through a
logarithm that nearly cancels between $u$ and $d$: the splitting fixes $w$ to a factor of two, no
better (R46).

The same rule fails for the pions: two circuits of three strands give 40~MeV at $\lone(6)$ or
508~MeV at $\lone(9)$, never 140. The pion is the closed ring at $r_e/2$, whose mode energy is
$2m_e c^2/\alpha = 140.05$~MeV ($\pi^\pm$: $+0.34\,\%$); this is the known coincidence
$m_\pi \approx 2m_e/\alpha$, an identification of the scale and not a derivation (R16).

\status{as in the table; the nucleon ring radius is the number the base does not fix}

%------------------------------------------------------------------------------
\subsection{The string tension and the nuclear force}\label{sec:tension}
%------------------------------------------------------------------------------

The energy per unit path of a matched line carrying the circulation quantum on one strand is
\[
u_{\rm path} = \frac{4\pi K q^2}{\lone(9)^2} = \frac{\pi\hbar c}{\lone(9)^2} = 939\ {\rm MeV/fm},
\qquad
\sqrt\sigma = \frac{3}{2\sqrt\pi}\,3^{2\pi}\,m_e c^2 = 430\ {\rm MeV},
\]
$\alpha$ and $Z_0$ cancelling between the pole charge and the matching (R39--R40,
\texttt{pitch\_search.py}, \texttt{circuit\_rule.py}). Lattice determinations put $\sqrt\sigma$ at
420--440~MeV \cite{necco}. The tension is the energy per path of the same object that weighs the
quark; the rule that assigns the quantum to one strand (a star arm is one circuit) is the one that
gives the electron its circulating half by the same formula, $\pi\hbar c/(3\lone(3)) = m_e c^2/2$.

The nuclear force is the pole-to-pole junction: a balanced junction of $k$ poles binds like
$\lfloor k/2\rfloor$ pairs (frustration leaves the odd pole inert), and a tri-junction of unit
poles at the nucleon size, $0.72$~fm, binds 2~MeV, the scale of the deuteron (2.22~MeV) and not of
confinement (R22, R24; \texttt{quark\_y.py}, \texttt{junction\_valence.py}). The manuscript's
statement that no confinement regime is reachable by winding is confirmed; the tension is not a
winding effect (R33--R39).

\status{Conditionally derived for $\sqrt\sigma$ (on the assignment of one quantum per star arm);
Derived for the $\lfloor k/2\rfloor$ law; Calibrated for the 2~MeV scale}

%------------------------------------------------------------------------------
\subsection{Lifetimes: the leak law and the weak wall}\label{sec:lifetime}
%------------------------------------------------------------------------------

Five readings of the author's ``breathing of the impedance'' were computed (R64--R69;
\texttt{breathing\_lifetime.py}, \texttt{resonance\_mismatch.py}, \texttt{current\_breathing.py},
\texttt{slowed\_reality.py}, \texttt{z\_field.py}, \texttt{stability\_law.py}). A self-similar
breathing gives $\tau \propto 1/m$ and puts the tau five orders too long: excluded. A resonance at
$Z \ne Z_0$ dies in fewer than a hundred turns: excluded (and this closes $Z_e = 0.73\,Z_0$ as a
circuit impedance). A relativistic slowing, global or internal, is either invisible or a velocity
tuned to $10^{-34}$: excluded. What the muon and tau impose is a leak per turn $\propto m^4$ with
the turn $T = 2\pi\hbar/mc^2$ (the Compton period, exact on the ladder), i.e. $\Gamma \propto m^5$;
an impedance even in the current, $Z(I) = Z_0[1 + (I/I_c)^2]$, the first correction a medium
without a preferred sense of circulation allows, gives that parity, but $I_c = 2.6\times10^8$~A is
not in the base.

The law assembled without adjustment,
\[
\tau(n) = T(n)\,\Big(\frac{M}{m(n)}\Big)^4\Big/ C(n),\qquad M = \frac{(96\pi^2)^{1/4}}{\sqrt{G_F}} = 1.625\ {\rm TeV},
\]
with $C(n)$ the number of lower rungs of the same charge class (colour counting three), gives the
electron stable ($C = 0$), the muon at $-0.4\,\%$, the tau at $+13\,\%$ by the bare count and
$+0.5\,\%$ with the standard QCD factor 1.2 on the quark channel, and the branching ratios
$e: 20\,\%$ (17.8), $\mu: 19.6\,\%$ (17.4), hadrons $60\,\%$ (64.8). The neutron, with the released
energy in place of the mass, comes out a factor 9 too long ($g_A$ and phase space, outside the
base). \textbf{$M$ is an identity}: with $\Gamma_\mu = G_F^2 m^5/192\pi^3$ and $T = 2\pi\hbar/m$,
$\Gamma T = (m/M)^4$ defines $M$; the law has the dimensional structure of the weak interaction and
does not derive $G_F$. The ladder has no rung at $M_W$ or $M_Z$ (nearest 76.8 and 104~GeV), and the
base has no neutral spin-1 excitation at $10^2$~GeV.

\status{Conditionally derived for the structure (Compton clock, $m^4$ leak by parity, channel
count); Reparametrized for $M$, which is $G_F$ under another name; Closed for a weak sector
inside the base}

%------------------------------------------------------------------------------
\subsection{Spin, the spinor bundle, and what statistics would need}\label{sec:spin}
%------------------------------------------------------------------------------

The manuscript's declared half-success was a magneton without spin $\tfrac12$, $g = 2$ or
statistics. The redraw built the following chain on the orientation of a closed ring (R74,
R78, R82--R85, R95--R98; \texttt{square\_odd\_mode.py}, \texttt{spin\_half\_rotor.py},
\texttt{core\_displacement.py}, \texttt{chern\_number.py}, \texttt{su2\_lift.py},
\texttt{mode\_separator.py}, \texttt{reorientation.py}), and it is reported with its locks:

\begin{itemize}[nosep]
\item A ring of three single strands has only $Z_3$, which has no element of order two and cannot
carry antiperiodicity; the square section has $Z_4$, and its odd mode, the $(1,2)/(2,1)$
Dirichlet doublet ($E$ irrep, $-1$ under $180^\circ$), with a half-twist of the section per
circuit, is antiperiodic. The half-twist $t = \tfrac12$ is posed. Curvature displaces the current
centroid 33.9~fm inward ($21.7\,\%$ of the section) under a stated constitutive condition (surface
current, constrained poloidal flux), so the dipolar pattern is not chosen; a rigid top is excluded
(a $j = 3/2$ would follow at $3mc^2$); the Chern number of the mode on the orientation sphere is 2
for the displaced core alone and 1 only with $t = \tfrac12$: it measures the half-twist, it does
not derive it.
\item Given $U(2\pi) = -1$ on the axis and the group law of rotations, quaternion tracking shows
that every $2\pi$ loop gives $-1$, the holonomy is $e^{i\Omega/2}$ and $c_1 = 1$ without Wigner's
theorem (R95): the SU(2) lift is derived from the axial sign and the group law.
\item Real analyzers built by projecting the field pattern never give $\cos^2(\theta/2)$
(deviation 0.25--0.44 after orthonormalization, up to $75\,\%$ of the norm outside the two channels,
R97); the half angle belongs to the transformation of the orientation, not to overlaps of fields.
\item With the only coupling the base has, $\mu_B$ in a field, the orientation precesses at
$\omega_L = eB/m$ and $\mathbf S\cdot\hat{\mathbf a}$ is conserved exactly; a classical
Stern--Gerlach (flat deflection distribution) and a damped alignment ($P(+) = 1$) are excluded by
experiment; two outcomes only, with conservation of $\mathbf S\cdot\hat{\mathbf a}$ in the mean,
give $p = \cos^2(\theta/2)$, the Born rule for one ring (R98). The discrete jump itself is not
derived.
\end{itemize}

The exchange sign between two rings, hence statistics, is not produced: the repository's own
analysis (\texttt{OPEN\_PROBLEMS.md} item 4) that a bosonic classical medium needs a superposition
of loop configurations with a sign rule stands. And after Section~\ref{sec:electron} the
carrier of this chain, the charged ring at $\lb$, is no longer the electron's electromagnetic
structure; what the chain establishes is a property of an orientable circulation, not of the
electron's charge.

\status{Conditionally derived for the SU(2) lift (on the posed half-twist) and for the Born rule
of one ring (on two outcomes); not supplied: $g = 2$ (an input after R103) and fermionic
statistics}

%------------------------------------------------------------------------------
\subsection{The electron: a representation excluded}\label{sec:electron}
%------------------------------------------------------------------------------

The redraw's electron (R53--R54) was a closed one-way ring of radius $\lb$ carrying the charge $e$,
half its rest energy circulating ($\pi\hbar c/2\pi\lb = m_e c^2/2$), half static in three
junctions; it gave $S = R\,E_{\rm circ}/c = \hbar/2$, $\mu = ecR/2 = \mu_B$, $g = 2$, no
radiation and no excited state. Four steps close it (\texttt{hydrogen\_ring.py},
\texttt{ring\_form\_factor.py}, \texttt{hydrogen\_form\_factor.py}, \texttt{moment\_without\_loop.py},
\texttt{electron\_is\_dirac.py}).

\begin{enumerate}[nosep]
\item \textbf{The atom of the base (R99, R101).} The base's atom is the Bohr orbit closed by de
Broglie's phase on the Lorentz-ether clock (R75). Its relativistic version is Sommerfeld's
\cite{sommerfeld}, and $E(n,k)$ coincides exactly with Dirac's $E(n,j)$ for $k = j + \tfrac12$
(checked to $10^{-9}$~eV on $1S$, $2S/2P_{1/2}$, $2P_{3/2}$, $3D_{5/2}$; fine structure 10.950~GHz
against 10.969 measured, the $0.18\,\%$ being QED). The base has the fine structure without the
ring. A rigid ring with a fixed axis would split the $1S$ by 263~GHz according to its orientation
(quadrupole in the proton's field) where only the 1.42~GHz hyperfine doublet exists: the
orientation must be a Bloch variable without quadrupole. The base's own orbits never come within
$a_0/2 = 68\,R$ of the proton, so the base's atom is blind to the ring; with the quantum-mechanical
$1S$ wave function imported, the ring shifts the $1S$ by $+205$ to $+232$~GHz against the 8.17~GHz
measured above Dirac.
\item \textbf{The form factor (R100), the author's criterion.} Starting from the charge as it
exists in the model and computing $F(\mathbf q) = e^{-1}\!\int\rho\,e^{i\mathbf q\cdot\mathbf r}d^3r$
and $\langle r^2\rangle = -6\,dF/dq^2|_0$, with PASS only if the existing theory gives
$F \approx 1$ despite the 386~fm ring: the three readings of the charge (thin ring at $\lb$; the
surface-current distribution of R84; $e/3$ at the three junctions) all give $F = \langle
j_0(q\rho)\rangle$ with $\langle r^2\rangle = 1.00$, $0.88$, $0.12\,\lb^2$; $F = 0.47$ at
1~MeV/$c$ and $\sim 10^{-4}$ at 1~GeV/$c$. Compton scattering at 1~MeV would be suppressed by
$|F|^2 = 0.18$ where Klein--Nishina \cite{kleinnishina} holds to the percent; elastic $e$--$p$ at
$Q^2 = 1$~GeV$^2$ by $10^{-8}$; LEP bounds the electron radius below $10^{-3}$~fm \cite{pdg}. The
orientation being a spin-$\tfrac12$ Bloch variable (item 1) closes the only direction in which a
ring is point-like ($\mathbf q$ along its axis); a snapshot reading (the probe sees the charge at
one point of the ring) leaves the elastic amplitude at $F$ and sends the rest into internal
excitations the electron does not have. The criterion fails.
\item \textbf{The Darwin term is not a size.} In the Breit frame $u^\dagger(p')u(p) = 2m$ exactly
at all $Q^2$: a point Dirac particle has $G_E = 1$ \cite{sachs}, and the $\tfrac34\lb^2$ of the
hydrogen Darwin term is the normalization factor $m/E = (1 + Q^2/4m^2)^{-1/2}$
\cite{foldy}, a kinematic quantity. The earlier prediction of the redraw (R89) that the ring should
have $r_{\rm rms} = (\sqrt3/2)\lb = 334$~fm to reproduce the Darwin term is withdrawn: the target
was never a charge radius.
\item \textbf{The magneton (R102).} A Bohr magneton carried by a charge loop, $\mu = qvR/2$ with
$q = e$ and $v \le c$, requires $R \ge \lb$ whatever the loop, while $F = 1$ to $1\,\%$ up to
1~GeV/$c$ requires $R < 0.05$~fm: $v/c \ge 8000$. The loop's magnetic form factor at $\lb$ is
0.72 at 1~MeV/$c$ and $10^{-6}$ at 1~GeV/$c$ against $G_M = 1$; the point Dirac particle has
$u^\dagger(p')\boldsymbol\alpha u(p) = i\,\boldsymbol\sigma\times\mathbf q$ exactly, $G_M = 1$, the
moment without a loop. The base has no other source of moment than a charge current.
\item \textbf{Identities (R103).} $S = \hbar/2$, $\mu_B$ and $g = 2$ follow only from the choice
$R = \lb$, $E_{\rm circ} = m/2$, $q = e$: with all the mass circulating $g = 1$, with the
Zitterbewegung radius $\lb/2$ and all the mass, $\mu_B/2$. The one proper output, the action
$\pi\hbar$ per turn, is the action per period of Dirac's Zitterbewegung (computed: $v = c$,
amplitude $\lb/2$, $\omega = 2mc^2/\hbar$, $E\,T = \pi\hbar$) \cite{schrodinger,hestenes}. The
static half in three junctions (R54) is the identity $3\hbar c/(\pi^2 D) = m_e c^2/2$ that
\emph{defines} $D = 6\lb/\pi^2$; its only content is $6/\pi^2 = 0.608$ against the manuscript's
$2\cosh\pi/37.1 = 0.625$, a $2.8\,\%$ coincidence.
\end{enumerate}

Decision: the electromagnetic electron is Dirac's. The base keeps what does not depend on where the
charge sits: the strand spectrum, the circulation energies (Sections~\ref{sec:ladder}--%
\ref{sec:tension}), the generation structure, the orientation chain. The ring at $\lb$ remains,
if anything, a neutral circulation of energy, the de Broglie clock, carrying no charge.

\status{Closed (charge loop, magneton loop); Identity (the ring's numbers); the decision is the
author's}

%------------------------------------------------------------------------------
\subsection{Three-dimensional stabilization}\label{sec:stabilization}
%------------------------------------------------------------------------------

No Derrick-evading term is derived. The stability the redraw uses is the conserved circulation
quantum per closed circuit (R47) together with the charge lock: a charged locked ring has no
free deformation that lowers its energy at fixed quantum, whereas a neutral locked defect has
$E = C/R$ with no minimum, no size, and pinned to the vacuum spacing weighs about 1~MeV and is
excluded as dark matter by nucleosynthesis (R88, R90, \texttt{dark\_defect.py}). The neutral,
transient mother loop of R87--R88 has no lock and decays into two locked daughters. This is a
postulate (one quantum per circuit) plus its consequences, not a stabilization mechanism in the
sense of \texttt{CONSTRAINTS.md}~\S3, and it says nothing about the medium's own elastic
coefficients.

\status{Postulated}

%------------------------------------------------------------------------------
\subsection{Exclusions recorded on the way}\label{sec:exclusions}
%------------------------------------------------------------------------------

\begin{itemize}[nosep]
\item Nucleons and quarks on the lepton ladder; masses as strand counts; ``one joint, one number''
(1.3~MeV at the nucleon, 4.6~MeV at the pion); additive masses calibrated on the electron; the
lepton ladder with a fourth rung.
\item Winding as a force mechanism; the flux tube in the base as written; a fixed linear charge
density; the vortex charge $5.85\,e$ as a circulating charge; free poles and standing waves (no
moment, radiation, harmonics); the $4$~MeV transverse mode (an artefact of the wave reading).
\item Lifetimes by self-similar breathing, by resonance at $Z \ne Z_0$, by relativistic slowing.
\item Dark matter as a neutral locked defect of the medium (no minimum; $\Gamma/H \sim 10^{20}$ at
nucleosynthesis if pinned); the medium's DQD in circulation ($5\times10^{33}$ times the observed
vacuum energy); an expanding medium ($g$ drift $10^{-10}$/yr against $3\times10^{-14}$) (R57, R90).
\item The base as a \emph{local} theory of Bell correlations: any definite-variable local model
gives $|S| \le 2$ (sign reading 2, Malus reading 0.94) against $2\sqrt2$ and the measured
$2.42 \pm 0.20$ \cite{hensen}; the shared quantum of a pair as a global constraint overshoots to
$S = 4$ without the Born rule; the medium's wave at $c_0 = c$ is $10^4$ times too slow to be the
channel; the joint state of two rings is not constructed (R91--R94).
\item The electron as an extended charge or an extended magneton (Section~\ref{sec:electron}).
\end{itemize}

%==============================================================================
\section{Epistemological status of the results}\label{sec:status}
%==============================================================================

\begin{table}[h]
\centering\small
\begin{tabular}{p{7.2cm}p{7.6cm}}
\toprule
level & claims filed at this level \\
\midrule
Derived & charge spectrum and colour from three strands (given the strand charges); the step of
four neutral strands and the exclusion of $n = 5$; three eigenvalues of a $C_3$ mass-root
operator; the $\{3,3,3\}$ partition; the strong part of $m_n - m_p$; the $\lfloor k/2\rfloor$
junction law; the SU(2) lift from the axial sign; that the ladder cannot be corrected into Koide;
that a loop magneton needs $R \ge \lb$. \\
\addlinespace
Conditionally derived & the quark scale and mass (on $2\pi$, 9); the nucleon mass and
$\Delta - N$ (on a ring radius); $m_n - m_p$ total (on the section and pole shape);
$\sqrt\sigma$ (on one quantum per star arm); Koide's $45^\circ$ (on the network reading of the
mass-root operator); the lifetime structure; the Born rule of one ring (on two outcomes); the
circulation quantum from a broken mother loop (on three inputs). \\
\addlinespace
Calibrated & the nuclear-force scale (2~MeV at 0.72~fm). \\
\addlinespace
Postulated & one quantum per closed circuit; the strand charges; the neutrino rule; the sharing
$2/3, 2/3, -1/3$; the half-twist $t = \tfrac12$; the phase $\varphi = 2/9$; the constitutive
condition of the surface current. \\
\addlinespace
Reparametrized & the exponent $2\pi$; the wall $M$ (an identity with $G_F$); the pion at $r_e/2$
(scale identification); the nucleon ring radius 0.67~fm. \\
\addlinespace
Closed & the extended electron (charge and magneton loops); a weak sector inside the base; dark
matter from the medium; the base as a local theory of Bell correlations; winding as confinement;
the mass-as-inductance law. \\
\bottomrule
\end{tabular}
\caption{The six-level ledger of this paper's claims. Nothing above is presented above its level.}
\end{table}

Two remarks on the ledger. First, every hadronic number rests on the exponent $2\pi$ and the count
9, which are inputs: the ratios are derived \emph{from} them, and a reader who does not grant them
should read the hadronic section as a parametrization with two structural inputs rather than a
derivation. Second, the coincidences recorded (proton radius as the envelope of three circuits,
$m_p c^2 = 4\hbar c/r_p$ to $0.05\,\%$, the nucleon ring at $\sqrt3$ circuit radii, $6/\pi^2$
against $2\cosh\pi/37.1$) are listed because they were found, not because they are claimed.

%==============================================================================
\section{Open problems}\label{sec:open}
%==============================================================================

\begin{enumerate}[nosep]
\item \textbf{The electron's charge and moment.} After R103 the base contains no source of a
point-like charge with a Bohr magneton. Either the base supplies a moment that is not a charge
current, or it adopts the Dirac structure and derives none of $S$, $\mu_B$, $g$. The neutral
circulation at $\lb$ (the de Broglie clock) is what remains of the ring.
\item \textbf{The phase $2/9$}, and a dynamics that gives the $45^\circ$ at the particle level
rather than the family level.
\item \textbf{The nucleon ring radius} (0.67~fm) and what the collective circulation of three
circuits turns on; the pion's pinning at $r_e/2$.
\item \textbf{The weak scale.} $G_F$, $M_W$, $M_Z$ and the mixing angle are not in the base; the
neutron's factor 9; the QCD factor 1.2 on the tau.
\item \textbf{Fermionic statistics}: the exchange sign between two rings; the discrete jump in a
measurement; the joint state of two rings (Bell).
\item \textbf{The exponent $2\pi$ and the count 9}, structural inputs.
\item \textbf{The section width} $w = 4\lb/\pi^2$, derived only through the three junctions of a
representation now closed; the splitting tests it to a factor of two.
\end{enumerate}

%==============================================================================
\section{Falsification}\label{sec:falsification}
%==============================================================================

One part of this sector has already been falsified by existing data: an electron whose charge or
magnetic moment is carried by a loop of any radius above $10^{-3}$~fm (Section~\ref{sec:electron}).
The paper reports this as its firmest result.

For what remains, the following measurements or computations would kill the corresponding claim:
\begin{itemize}[nosep]
\item a lattice determination of the isospin-breaking QCD part of $m_n - m_p$ that departs from
$2.47$~MeV by more than its present $\pm 0.29$~MeV, or of $\sqrt\sigma$ outside 420--440~MeV by more
than the $2\,\%$ the exponent allows (the formulae $(2\alpha/3)3^{2\pi} m_e c^2$ and
$(3/2\sqrt\pi)3^{2\pi} m_e c^2$ have no free number);
\item a charged lepton of any mass below the ladder's next rung would not falsify the $C_3$
structure, but a violation of Koide's relation beyond $10^{-5}$ by an improved $\tau$ mass would
remove the $45^\circ$ that the $g = 2$ pairing explains;
\item a $\Delta$--$N$ gap or a nucleon magnetic-moment ratio that no single ring radius fits within
$15\,\%$ and $3\,\%$ respectively;
\item a tau lifetime or branching ratios that the channel count $C(n)$ cannot reach within the
$13\,\%$ the bare count leaves, once the QCD factor is computed rather than borrowed.
\end{itemize}
Nothing in the generational sector predicts a new particle; the sector's exposure is entirely
through the precision of the ratios above.

%==============================================================================
\section*{Acknowledgments and Methodological Disclosure}\label{sec:disclosure}
%==============================================================================

The main manuscript discloses its human--AI collaboration and states its limitation: both
numerical implementations originate from the same collaboration, and independent third-party
verification has not been performed. This paper involves a second AI system, and the case that
applies is the first of the two the repository distinguishes: \textbf{the system was given the
author's statements and the repository, including the manuscript's constraints, and asked to
compute their consequences.} The prior is shared. This is not independent verification, and it is
not described as such anywhere in this paper.

Roles. The author supplied the constitutive statements (R1--R8 and the later ones quoted in the
record), the direction at each step, and the review comments that imposed the downgrades recorded
as ``corrections'' in \texttt{matter/redraw/BASE.md} (R68 from derived to identity; R86 from
derived to conditional; R92--R93, R96--R97 rewordings; the criterion of R100). The second system
(Claude, Anthropic, used through Claude Code sessions in September 2026; the model designation is
in the session record and is to be filled in by the author) wrote the 98 scripts, ran them,
wrote the record and this draft. Every number in this paper is produced by a script in
\texttt{matter/redraw/} that prints a \texttt{RESULT} line; one script,
\texttt{ring\_form\_factor.py}, is left with a failing check on purpose, the failure being the
author's criterion and the result of R100.

%==============================================================================
\begin{thebibliography}{99}
\bibitem{pdg} S. Navas \textit{et al.} (Particle Data Group), \textit{Review of Particle Physics}, Phys. Rev. D \textbf{110}, 030001 (2024).
\bibitem{dsmodel} R. St-Pierre, \textit{Dipolar String Model: an Exploratory Framework for the Vacuum from Transmission-Line Electrodynamics}, V2.10 (2026). Concept DOI: \href{https://doi.org/10.5281/zenodo.21196098}{10.5281/zenodo.21196098}.
\bibitem{koide} Y. Koide, Lett. Nuovo Cimento \textbf{34}, 201 (1982); Phys. Lett. B \textbf{120}, 161 (1983).
\bibitem{foot} R. Foot, \textit{A note on Koide's lepton mass relation}, arXiv:hep-ph/9402242 (1994).
\bibitem{brannen} C. A. Brannen, \textit{The lepton masses} (2006), unpublished; the phase $2/9$ in the parametrization $\sqrt{m_k} \propto 1 + \sqrt2\cos(2/9 + 2\pi k/3)$.
\bibitem{bilsonthompson} S. O. Bilson-Thompson, \textit{A topological model of composite preons}, arXiv:hep-ph/0503213 (2005).
\bibitem{bmw} Sz. Borsanyi \textit{et al.} (BMW Collaboration), \textit{Ab initio calculation of the neutron--proton mass difference}, Science \textbf{347}, 1452 (2015).
\bibitem{necco} S. Necco and R. Sommer, Nucl. Phys. B \textbf{622}, 328 (2002); the range 420--440~MeV quoted here spans the quenched and unquenched determinations of $\sqrt\sigma$.
\bibitem{katrin} M. Aker \textit{et al.} (KATRIN Collaboration), \textit{Direct neutrino-mass measurement based on 259 days of KATRIN data}, Science \textbf{388}, 180 (2025).
\bibitem{hensen} B. Hensen \textit{et al.}, \textit{Loophole-free Bell inequality violation using electron spins separated by 1.3 kilometres}, Nature \textbf{526}, 682 (2015).
\bibitem{sommerfeld} A. Sommerfeld, Ann. Phys. (Leipzig) \textbf{51}, 1 (1916).
\bibitem{kleinnishina} O. Klein and Y. Nishina, Z. Phys. \textbf{52}, 853 (1929).
\bibitem{sachs} R. G. Sachs, Phys. Rev. \textbf{126}, 2256 (1962).
\bibitem{foldy} L. L. Foldy and S. A. Wouthuysen, Phys. Rev. \textbf{78}, 29 (1950).
\bibitem{schrodinger} E. Schr\"odinger, Sitzungsber. Preuss. Akad. Wiss., Phys.-Math. Kl. \textbf{24}, 418 (1930).
\bibitem{hestenes} D. Hestenes, \textit{Zitterbewegung in quantum mechanics}, Found. Phys. \textbf{40}, 1 (2010).
\end{thebibliography}

\end{document}
